% ===========================================================================
%  Artículo Académico – Sistema Wearable de Monitoreo de Actividades
%  de la Vida Diaria basado en ESP32 con ESP-NOW, IMU y Registro en SD
%  Universidad del Cauca — Departamento de Ingeniería Electrónica
%  y Telecomunicaciones
%  Autor: Carlos Daniel Cuellar Antury
%  Formato: doble columna, 10 pt, Times New Roman, márgenes 2.5 cm, 6 páginas
% ===========================================================================
\documentclass[10pt,twocolumn,a4paper]{article}

% --- Paquetes ---------------------------------------------------------------
\usepackage[utf8]{inputenc}
\usepackage[T1]{fontenc}
\usepackage[spanish,es-tabla]{babel}
\usepackage{mathptmx}
\usepackage[top=2.5cm,bottom=2.5cm,left=2.5cm,right=2.5cm,
            columnsep=0.6cm]{geometry}
\usepackage{graphicx}
\usepackage{booktabs}
\usepackage{array}
\usepackage{listings}
\usepackage{xcolor}
\usepackage{hyperref}
\usepackage{caption}
\usepackage{float}
\usepackage{amsmath}
\usepackage{fancyhdr}
\usepackage{titlesec}
\usepackage{enumitem}
\usepackage{microtype}
\usepackage{url}

% --- Encabezado y pie -------------------------------------------------------
\pagestyle{fancy}
\fancyhf{}
\fancyhead[L]{\tiny Univ. del Cauca — Ing. Electrónica y Telecomunicaciones}
\fancyhead[R]{\tiny Cuellar Antury, C.D. (2024)}
\fancyfoot[C]{\small\thepage}
\renewcommand{\headrulewidth}{0.4pt}
\setlength{\headheight}{14pt}

% --- Formato de secciones ---------------------------------------------------
\titleformat{\section}{\normalfont\bfseries\small}{\thesection.}{0.4em}{}
\titleformat{\subsection}{\normalfont\bfseries\footnotesize}{\thesubsection.}{0.4em}{}
\titleformat{\subsubsection}{\normalfont\itshape\footnotesize}{\thesubsubsection.}{0.4em}{}
\titlespacing{\section}{0pt}{5pt}{2pt}
\titlespacing{\subsection}{0pt}{3pt}{1pt}
\titlespacing{\subsubsection}{0pt}{2pt}{1pt}

% --- Listas compactas -------------------------------------------------------
\setlist{itemsep=1pt,topsep=2pt,parsep=0pt,partopsep=0pt}

% --- Colores para código ----------------------------------------------------
\definecolor{codegray}{rgb}{0.96,0.96,0.96}
\definecolor{codeblue}{rgb}{0.0,0.0,0.55}
\definecolor{codegreen}{rgb}{0.0,0.45,0.0}
\definecolor{codepurple}{rgb}{0.45,0.0,0.45}

\lstdefinestyle{csrc}{
  language=C,
  backgroundcolor=\color{codegray},
  basicstyle=\ttfamily\scriptsize,
  keywordstyle=\color{codeblue}\bfseries,
  commentstyle=\color{codegreen}\itshape,
  stringstyle=\color{codepurple},
  numberstyle=\tiny\color{gray},
  numbers=left,
  stepnumber=1,
  numbersep=4pt,
  breaklines=true,
  showspaces=false,
  showstringspaces=false,
  frame=single,
  rulecolor=\color{black!20},
  tabsize=2,
  captionpos=b,
  aboveskip=4pt,
  belowskip=2pt,
  literate=
    {á}{{\'a}}1 {é}{{\'e}}1 {í}{{\'i}}1 {ó}{{\'o}}1 {ú}{{\'u}}1
    {Á}{{\'A}}1 {É}{{\'E}}1 {Í}{{\'I}}1 {Ó}{{\'O}}1 {Ú}{{\'U}}1
    {ñ}{{\~n}}1 {Ñ}{{\~N}}1
}

\hypersetup{colorlinks=true,linkcolor=blue,urlcolor=blue,citecolor=blue,
            breaklinks=true}
\captionsetup{font=scriptsize,labelfont=bf,skip=3pt}

% ===========================================================================
\begin{document}

% ---------------------------------------------------------------------------
%  PORTADA — TÍTULO, AUTOR, AFILIACIÓN, ABSTRACT (una columna)
% ---------------------------------------------------------------------------
\twocolumn[{%
\begin{center}
  {\large\bfseries Sistema Wearable de Monitoreo de Actividades de la Vida\\[2pt]
  Diaria basado en ESP32: Arquitectura Distribuida con ESP-NOW,\\[2pt]
  IMU, Pantalla Táctil y Registro en Tarjeta SD}\\[6pt]
  {\normalsize Carlos Daniel Cuellar Antury}\\[2pt]
  {\small Departamento de Ingeniería Electrónica y Telecomunicaciones,
  Universidad del Cauca, Popayán, Colombia}\\[2pt]
  {\footnotesize Repositorio: \url{https://github.com/Cdcuellar/Sistemas-embebidos---Actividades-de-la-vida-diaria}}\\[6pt]
  \rule{\linewidth}{0.5pt}\\[4pt]
\end{center}
\begin{small}
\noindent\textbf{Resumen} — Este artículo describe el diseño e implementación de un
sistema embebido distribuido para el monitoreo de Actividades de la Vida Diaria
(AVD), desarrollado como proyecto académico en la Universidad del Cauca. El
sistema adopta una arquitectura de tres nodos ESP32 interconectados mediante el
protocolo ESP-NOW: un nodo maestro que integra un sensor de frecuencia cardíaca
(ADC), reloj de tiempo real DS3231, almacenamiento en tarjeta microSD y
comunicación inalámbrica; un nodo esclavo de movimiento con sensor inercial
MPU6050 (acelerómetro y giroscopio 6-DOF); y un nodo esclavo con pantalla
circular táctil GC9A01 de 1,28'' y sensor QMI8658. Adicionalmente, se documenta
una herramienta de diagnóstico que identifica el registro \texttt{WHO\_AM\_I}
del MPU6050 y una utilidad de obtención de dirección MAC. El firmware se
desarrolló con ESP-IDF~v6.0, FreeRTOS y gestión de energía con escalado dinámico
de frecuencia. Los resultados muestran transmisión estable a 100~ms (IMU) y
50~ms (touch/IMU-S3), con registro CSV local y visualización en tiempo real.

\smallskip
\noindent\textbf{Palabras clave} — Wearable, AVD, ESP32, ESP-NOW, MPU6050,
QMI8658, FreeRTOS, ESP-IDF, tarjeta SD, I2C, pantalla táctil, frecuencia
cardíaca.

\smallskip
\noindent\textbf{Abstract} — This paper describes the design and implementation
of a distributed embedded system for monitoring Activities of Daily Living (ADL),
developed as an academic project at Universidad del Cauca. The system adopts a
three-node ESP32 architecture interconnected via the ESP-NOW protocol: a master
node integrating a heart-rate sensor (ADC), DS3231 real-time clock, microSD
storage, and wireless communication; a motion slave node with an MPU6050
inertial sensor (6-DOF accelerometer and gyroscope); and a slave node with a
1.28'' GC9A01 circular touch display and QMI8658 sensor. Additionally, a
diagnostic tool that identifies the MPU6050 \texttt{WHO\_AM\_I} register and a
MAC address utility are documented. Firmware was developed with ESP-IDF~v6.0,
FreeRTOS, and dynamic frequency scaling power management. Results show stable
transmission at 100~ms (IMU node) and 50~ms (touch/IMU-S3 node), with local CSV
logging and real-time display.

\smallskip
\noindent\textbf{Keywords} — Wearable, ADL, ESP32, ESP-NOW, MPU6050, QMI8658,
FreeRTOS, ESP-IDF, SD card, I2C, touch display, heart rate.
\end{small}
\vspace{4pt}\rule{\linewidth}{0.5pt}\vspace{4pt}
}]

% ---------------------------------------------------------------------------
%  1. INTRODUCCIÓN
% ---------------------------------------------------------------------------
\section{Introducción}

El envejecimiento poblacional, el aumento de las enfermedades crónicas y
el interés creciente por la telemedicina han impulsado el desarrollo de
sistemas wearables capaces de monitorear de forma continua y no intrusiva
las Actividades de la Vida Diaria (AVD) de los pacientes~\cite{chen2012}.
Estos sistemas registran señales biométricas y biomecánicas —frecuencia
cardíaca, postura, movimiento— para apoyar diagnósticos clínicos,
programas de rehabilitación y asistencia a personas mayores.

En este contexto, la familia de microcontroladores ESP32 (Espressif
Systems) presenta características sobresalientes para aplicaciones
wearables: bajo consumo, conectividad inalámbrica nativa (WiFi +
Bluetooth), múltiples interfaces de periféricos y un ecosistema de
desarrollo maduro basado en ESP-IDF~\cite{espressif2023}. El protocolo
ESP-NOW, en particular, permite la transmisión directa punto a punto de
tramas de hasta 250~bytes entre dispositivos ESP32 sin necesidad de un
router intermedio, con latencias del orden de los milisegundos y un
consumo energético significativamente menor al de WiFi
convencional~\cite{espnow_docs}.

El presente artículo documenta el diseño completo de un sistema wearable
distribuido para monitoreo de AVD, publicado como proyecto de código
abierto en GitHub~\cite{repo_github}. El sistema comprende cuatro
submódulos funcionales: (1)~un nodo maestro encargado de centralizar
datos, sincronizar la hora con un RTC y registrar en SD; (2)~un nodo
esclavo de movimiento basado en MPU6050; (3)~un nodo esclavo con
pantalla circular táctil y sensor QMI8658; y (4)~herramientas de apoyo
para diagnóstico de hardware.

\subsection{Objetivos}

\begin{itemize}
  \item Diseñar una arquitectura distribuida de tres nodos ESP32
        interconectados por ESP-NOW para captura y centralización de
        datos de AVD.
  \item Implementar firmware con FreeRTOS, gestión de energía y
        comunicación robusta en cada nodo.
  \item Documentar el código fuente, esquemas de conexión y resultados
        para su reproducibilidad pública en GitHub.
\end{itemize}

% ---------------------------------------------------------------------------
%  2. ESTADO DEL ARTE
% ---------------------------------------------------------------------------
\section{Estado del Arte}

El reconocimiento de AVD mediante sistemas inerciales es un campo
ampliamente estudiado. Anguita \textit{et al.}~\cite{anguita2012}
publicaron un \textit{dataset} de referencia con acelerómetro de
teléfono inteligente y 30 sujetos, logrando tasas de clasificación
$>$96\% con SVM para seis clases (caminar, subir escaleras, sentarse,
acostarse, estar de pie, tumbarse). Chen \textit{et al.}~\cite{chen2012}
revisaron sistemáticamente más de 200 trabajos en reconocimiento de
actividad basado en sensores, identificando IMUs de bajo coste, fusión
sensorial y clasificación automática como las tendencias dominantes
para aplicaciones de monitoreo de salud ambiental y asistida.

El ESP32 se ha posicionado como plataforma de referencia para sistemas
wearables académicos e industriales de bajo coste. Kolban~\cite{kolban2017}
documenta en detalle la arquitectura dual-core, el sistema de gestión
de energía con sus modos activo, modem-sleep, light-sleep, deep-sleep y
hibernación, y la pila de red WiFi/BT. El protocolo ESP-NOW permite la
transmisión directa punto a punto de tramas de hasta 250~bytes entre
ESP32 sin router intermedio, con latencias de 2--5~ms y un consumo
significativamente menor al de WiFi convencional~\cite{espnow_docs}.

La combinación de ESP-NOW con FreeRTOS y gestión dinámica de frecuencia
(DFS) ha sido reportada como estrategia efectiva para extender la vida
útil de la batería sin comprometer la tasa de muestreo en aplicaciones
de redes de área corporal (BAN). El presente trabajo adopta esta
arquitectura e implementa todos sus componentes con el nivel de detalle
requerido para su reproducción pública mediante un repositorio GitHub
de acceso abierto~\cite{repo_github}.

Respecto al sensor MPU6050~\cite{invensense2013}, su uso masivo en el
mercado de componentes ha generado un problema documentado de clones
con identificadores \texttt{WHO\_AM\_I} distintos al estándar
\texttt{0x68}. Este trabajo aborda dicha problemática mediante una
herramienta de diagnóstico dedicada que se integra al flujo de puesta
en marcha del hardware, contribuyendo a la trazabilidad experimental.

La pantalla circular GC9A01 de 1,28'' ha emergido como opción atractiva
para wearables compactos por su bajo consumo, resolución 240$\times$240~px
y soporte para renderizado eficiente por SPI a 40~MHz. La combinación
con el ESP32-S3 y la IMU QMI8658 en módulos integrados abre nuevas
posibilidades para interfaces de usuario en monitoreo biomédico.

% ---------------------------------------------------------------------------
%  3. ARQUITECTURA DEL SISTEMA
% ---------------------------------------------------------------------------
\section{Arquitectura del Sistema}
\label{sec:arq}

El sistema adopta una topología estrella con un nodo maestro y dos
nodos esclavos. La Figura~\ref{fig:arq} ilustra la arquitectura
general.

\begin{figure}[H]
\centering
\scriptsize
\begin{verbatim}
 +---------------------------+
 |   NODO MAESTRO (ESP32)    |
 |  ADC: Pulso (GPIO34)      |
 |  I2C: DS3231 RTC          |
 |  SPI: microSD (datos.csv) |
 |  ESP-NOW: Recepción       |
 +----+-----------+----------+
      |           |
   ESP-NOW     ESP-NOW
      |           |
 +----+---+  +----+------+
 | NODO 2 |  |  NODO 3   |
 | MPU6050|  | QMI8658   |
 | I2C    |  | GC9A01    |
 | 100 ms |  | CST816S   |
 +--------+  | 50 ms     |
             +-----------+
\end{verbatim}
\caption{Arquitectura distribuida del sistema wearable AVD}
\label{fig:arq}
\end{figure}

La Tabla~\ref{tab:nodos} resume las características principales de
cada nodo.

\begin{table}[H]
\centering
\caption{Resumen de nodos del sistema}
\label{tab:nodos}
\scriptsize
\begin{tabular}{p{1.2cm}p{1.4cm}p{3.1cm}}
\toprule
\textbf{Nodo} & \textbf{MCU} & \textbf{Función principal} \\
\midrule
Maestro & ESP32 & Centraliza datos, RTC, SD, pulso \\
Esclavo 2 & ESP32 & IMU MPU6050, roll/pitch, ESP-NOW \\
Esclavo 3 & ESP32-S3 & IMU QMI8658, pantalla táctil, ESP-NOW \\
Herr. & ESP32 & Diagnóstico WHO\_AM\_I (Arduino IDE) \\
MAC & ESP32 & Obtención de dirección MAC WiFi \\
\bottomrule
\end{tabular}
\end{table}

\subsection{Protocolo de Comunicación ESP-NOW}

ESP-NOW opera sobre la capa física 802.11 en modo ad-hoc sin
asociación a un punto de acceso~\cite{espnow_docs}. Los paquetes
tienen una carga útil máxima de 250~bytes y viajan con confirmación
de entrega (ACK de capa 2). En este sistema, todos los nodos utilizan
la misma estructura de 32~bytes (8~floats) para los datos IMU:

\begin{lstlisting}[style=csrc, numbers=none,
  caption={Estructura de datos transmitida por ESP-NOW (32 bytes)},
  label=lst:struct]
typedef struct {
    float roll, pitch;       // Angulos (grados)
    float acc_x, acc_y, acc_z; // Aceleracion (g)
    float gyr_x, gyr_y, gyr_z; // Giro (deg/s)
} datos_imu_t; // 8 x 4 bytes = 32 bytes exactos
\end{lstlisting}

La consistencia del tamaño se verifica en tiempo de compilación con
\texttt{\_Static\_assert(sizeof(...) == 32, ...)}, garantizando la
compatibilidad entre nodos incluso frente a cambios futuros.

% ---------------------------------------------------------------------------
%  4. DISEÑO E IMPLEMENTACIÓN POR NODO
% ---------------------------------------------------------------------------
\section{Diseño e Implementación}
\label{sec:impl}

\subsection{Nodo Maestro (esp32\_1\_hr\_rstc\_sd\_master)}

El nodo maestro es el controlador central del sistema. Ejecuta múltiples
tareas FreeRTOS concurrentes bajo un único mutex para proteger el
contenedor de datos compartidos (\texttt{log\_completo\_t}).

\subsubsection{Sensor de Frecuencia Cardíaca}

La frecuencia cardíaca se adquiere localmente mediante el ADC del ESP32
en el canal 6 (GPIO~34). Se implementa un filtro en dos etapas:
\begin{enumerate}
  \item Ráfaga de 8 muestras con media recortada (se descartan máximo y
        mínimo).
  \item Filtro de media exponencial (EMA con $\alpha=0{,}85$) sobre la
        señal resultante, seguido de eliminación de componente DC con
        una segunda EMA ($\alpha=0{,}98$) como línea de base.
\end{enumerate}

Esta cadena de procesamiento reduce el ruido eléctrico inherente a la
lectura analógica sobre el bus de alimentación del ESP32.

\subsubsection{Reloj de Tiempo Real DS3231}

El DS3231 se comunica vía I2C (GPIO~21 SDA, GPIO~22 SCL, dirección
\texttt{0x68}) a 100~kHz. Los registros de tiempo se leen en formato
BCD y se convierten a \texttt{struct~tm}. El firmware permite
sincronización de hora desde consola mediante el comando \texttt{SETTIME}
por puerto serie.

\subsubsection{Registro en microSD}

Los datos se acumulan en un buffer RAM de 60 filas y se vuelcan cada
minuto al archivo \texttt{/sdcard/datos.csv} mediante el sistema de
archivos FAT (ESP-VFS-FAT). La interfaz SPI se configura en el slot
SPI3 (GPIO~23 MOSI, GPIO~19 MISO, GPIO~18 SCK, GPIO~5 CS). Esta
estrategia reduce el número de operaciones de escritura en flash y
mejora la vida útil de la tarjeta.

\subsubsection{Gestión de Energía}

El maestro implementa escalado dinámico de frecuencia (DFS):
CPU a 240~MHz en modo activo, 80~MHz en modo normal y 10~MHz en modo
inactivo. Adicionalmente se habilita \textit{WiFi Power Save}
(\texttt{WIFI\_PS\_MIN\_MODEM}, ahorro $\sim$20~mA) y \textit{Light
Sleep} con fallback si el hardware no lo soporta. El log en consola se
reduce a cada 10~s para minimizar el overhead de UART.

\subsection{Nodo Esclavo 2 — Movimiento (esp32\_2)}

El nodo de movimiento adquiere datos del sensor MPU6050 cada 100~ms y
los envía al maestro cada 500~ms mediante ESP-NOW. Utiliza dos tareas
FreeRTOS protegidas por un mutex binario:

\begin{itemize}
  \item \textbf{Tarea de muestreo}: lee los 14~bytes del MPU6050
        (acelerómetro + giroscopio + temperatura) en una única
        transacción I2C, aplica los offsets de calibración y calcula
        roll/pitch.
  \item \textbf{Tarea de envío}: copia atómicamente los datos con el
        mutex y los transmite por ESP-NOW.
\end{itemize}

\subsubsection{Driver MPU6050 (mpu6050.c/h)}

Se implementó un driver C completo y documentado con Doxygen para el
MPU6050 usando el nuevo driver I2C master de ESP-IDF~v5+
(\texttt{driver/i2c\_master.h}). Sus funciones principales son:

\begin{table}[H]
\centering
\caption{API del driver MPU6050}
\label{tab:api_mpu}
\scriptsize
\begin{tabular}{p{2.8cm}p{3.9cm}}
\toprule
\textbf{Función} & \textbf{Descripción} \\
\midrule
\texttt{mpu6050\_init()} & Inicializa I2C, verifica WHO\_AM\_I, configura rangos \\
\texttt{mpu6050\_calibrate()} & Promedia N muestras, calcula y almacena offsets \\
\texttt{mpu6050\_read\_accel\_raw()} & Lee 3 ejes del acelerómetro (6 bytes) \\
\texttt{mpu6050\_read\_all\_raw()} & Lee los 14 bytes completos (accel + temp + gyro) \\
\bottomrule
\end{tabular}
\end{table}

La función \texttt{mpu6050\_init()} acepta los valores \texttt{0x68},
\texttt{0x69} y \texttt{0x72} en el registro \texttt{WHO\_AM\_I}
(dirección \texttt{0x75}), cubriendo tanto el MPU6050 original como
variantes de silicio detectadas durante el desarrollo.

\subsubsection{Calibración del MPU6050}

La calibración del sensor se realiza automáticamente al arrancar el
Nodo~2, mediante la función \texttt{mpu6050\_calibrate()}, que toma
100~muestras con el sensor estático y calcula los offsets de cada eje:

\begin{equation}
  \text{offset}_i = \frac{1}{N}\sum_{k=1}^{N} \text{raw}_{i,k} \qquad
  (i \in \{a_x, a_y, a_z, g_x, g_y, g_z\})
  \label{eq:calib}
\end{equation}

El eje $a_z$ se corrige restando el equivalente a 1~g en LSB
(16384 para rango $\pm$2~g) para que el reposo sobre superficie plana
reporte $a_z = 1\text{ g}$ y no $0\text{ g}$.

Los offsets quedan almacenados en la estructura \texttt{mpu6050\_t}
y se aplican automáticamente en cada llamada a
\texttt{mpu6050\_read\_all\_raw()}, sin necesidad de post-procesamiento
en la aplicación.

\subsubsection{Robustez en la Comunicación I2C}

El driver implementa reintentos automáticos con hasta 5 intentos y
5~ms de espera entre cada uno para lecturas de registros. Para la
lectura del \texttt{WHO\_AM\_I} en la inicialización se permite hasta
20~reintentos, dado que algunos clones del sensor requieren más tiempo
para responder tras el reset. Esta estrategia mejora significativamente
la tasa de éxito en la inicialización del sistema con hardware de
terceros.

\begin{equation}
  \text{Roll} = \arctan\!\left(\frac{a_y}{a_z}\right)\cdot\frac{180°}{\pi}
  \label{eq:roll}
\end{equation}
\begin{equation}
  \text{Pitch} = \arctan\!\left(\frac{-a_x}{\sqrt{a_y^2+a_z^2}}\right)\cdot\frac{180°}{\pi}
  \label{eq:pitch}
\end{equation}

Estas fórmulas proporcionan ángulos estables en condiciones cuasi-estáticas
(sin vibración) con rango $\pm90°$.

\subsection{Nodo Esclavo 3 — Touch/Display (esp\_32\_s3\_touch\_1.28)}

El tercer nodo usa un ESP32-S3 con pantalla circular GC9A01 de 1,28''
(240$\times$240 px, SPI a 40~MHz) y sensor táctil capacitivo CST816S.
Su función dual es: (a)~visualizar un nivel de burbuja animado en tiempo
real; y (b)~transmitir datos de movimiento al maestro por ESP-NOW a 50~ms
(20~Hz). El nodo también mide el voltaje de batería mediante un divisor
resistivo en el ADC.

\subsubsection{Sensor IMU QMI8658}

El QMI8658 es la IMU interna del módulo ESP32-S3 Wearable 1.28'', con
interfaz I2C en GPIO~6 (SDA) y GPIO~7 (SCL) a 400~kHz (dirección
\texttt{0x6B}). Sus rangos configurados son $\pm$8~g para el
acelerómetro (4096~LSB/g) y $\pm$512~°/s para el giroscopio (64~LSB/(°/s)).

\subsubsection{Filtro Complementario}

Para estimar los ángulos con mayor precisión que las ecuaciones~\ref{eq:roll}
y~\ref{eq:pitch}, se aplica un filtro complementario:

\begin{equation}
  \theta_k = \alpha_g(\theta_{k-1} + \omega\,\Delta t) + \alpha_a\,\theta_{\text{acc}}
  \label{eq:comp}
\end{equation}

donde $\alpha_g=0{,}98$ (peso del giroscopio) y $\alpha_a=0{,}02$ (peso
del acelerómetro), con período $\Delta t=50$~ms. Este filtro corrige la
deriva del giroscopio usando la inclinación de largo plazo del
acelerómetro.

\subsubsection{Nivel de Burbuja Visual}

La posición de la burbuja en pantalla se mapea directamente a los
ángulos roll/pitch según:
$x_b = 120 + \text{roll}\cdot k$,
$y_b = 120 - \text{pitch}\cdot k$,
donde $k$ es un factor de escala de píxeles por grado. La pantalla
se apaga automáticamente tras 10~s sin interacción táctil para
ahorrar energía.

\subsection{FreeRTOS y Sincronización entre Tareas}

La arquitectura de software de cada nodo se basa en FreeRTOS con
múltiples tareas concurrentes protegidas por mutexes. La
Tabla~\ref{tab:tareas} resume las tareas de cada nodo.

\begin{table}[H]
\centering
\caption{Tareas FreeRTOS por nodo}
\label{tab:tareas}
\scriptsize
\begin{tabular}{p{1.5cm}p{2.0cm}p{2.7cm}}
\toprule
\textbf{Nodo} & \textbf{Tarea} & \textbf{Función} \\
\midrule
Maestro & task\_rx\_espnow & Recepción paquetes ESP-NOW \\
Maestro & task\_rtc & Lectura periódica del DS3231 \\
Maestro & task\_adc\_pulso & Adquisición pulso + filtrado \\
Maestro & task\_sd\_log & Volcado CSV a microSD \\
Nodo 2 & tarea\_muestreo & Lectura MPU6050 cada 100 ms \\
Nodo 2 & tarea\_envio & Envío ESP-NOW cada 500 ms \\
Nodo 3 & tarea\_imu & Lectura QMI8658 + filtro comp. \\
Nodo 3 & tarea\_espnow & Envío paquete cada 50 ms \\
Nodo 3 & tarea\_display & Renderizado pantalla GC9A01 \\
\bottomrule
\end{tabular}
\end{table}

El patrón de sincronización es consistente en todos los nodos: la tarea
productora (lectura del sensor) adquiere el mutex, actualiza el buffer
compartido y libera el mutex; la tarea consumidora (envío o registro)
también adquiere el mutex, copia los datos y los procesa de forma
independiente. Este desacoplamiento evita que una operación lenta (p.ej.
escritura en SD o transmisión ESP-NOW) bloquee la adquisición de datos.

\subsection{Estructura del Repositorio GitHub}

La Tabla~\ref{tab:repo} documenta la estructura de carpetas del
repositorio y la función de cada módulo.

\begin{table}[H]
\centering
\caption{Estructura del repositorio GitHub}
\label{tab:repo}
\scriptsize
\begin{tabular}{p{3.0cm}p{3.7cm}}
\toprule
\textbf{Directorio} & \textbf{Contenido} \\
\midrule
\texttt{esp32\_1\_hr\_rstc\_sd\_master/} & Nodo maestro: RTC, SD, pulso, ESP-NOW \\
\texttt{esp32\_2/} & Nodo IMU: MPU6050, driver C, ESP-NOW \\
\texttt{esp\_32\_s3\_touch\_1.28/} & Nodo touch: QMI8658, GC9A01, CST816S \\
\texttt{Direccion mac/} & Utilidad de lectura de MAC WiFi \\
\texttt{herramienta\_who\_i\_am\_*/} & Diagnóstico WHO\_AM\_I (Arduino IDE) \\
\bottomrule
\end{tabular}
\end{table}

Cada directorio de proyecto ESP-IDF contiene \texttt{CMakeLists.txt},
\texttt{Doxyfile} para generación de documentación automática,
\texttt{.devcontainer/} para entorno Docker reproducible,
y \texttt{.vscode/} con configuración de depuración remota por OpenOCD.
\label{sec:who}

Durante la fase de integración del Nodo~2, se detectó que la biblioteca
de producción rechazaba el MPU6050 por discrepancia en el registro
\texttt{WHO\_AM\_I}. Se desarrolló una herramienta mínima en Arduino
IDE que lee directamente el registro \texttt{0x75} vía I2C y reporta
el valor real del silicio por el monitor serie a 115\,200~baudios. El
valor obtenido fue \texttt{0x72}, que difiere del \texttt{0x68}
estándar del datasheet. Esta información se incorporó al driver de
producción \texttt{mpu6050.c} como condición de aceptación.

\subsection{Utilidad de Dirección MAC (Direccion\_mac)}

Para configurar correctamente la tabla de peers de ESP-NOW, se desarrolló
una utilidad independiente en ESP-IDF que inicializa WiFi en modo STA y
reporta la dirección MAC de 6~bytes del ESP32 por el puerto serie:

\begin{lstlisting}[style=csrc, numbers=none,
  caption={Lectura de MAC WiFi en ESP-IDF},
  label=lst:mac]
uint8_t mac[6];
esp_wifi_get_mac(WIFI_IF_STA, mac);
printf("MAC: %02x:%02x:%02x:%02x:%02x:%02x\n",
       mac[0],mac[1],mac[2],
       mac[3],mac[4],mac[5]);
\end{lstlisting}

Las MACs obtenidas se usaron para configurar los arrays
\texttt{mac\_esclavo\_mpu} y \texttt{mac\_esclavo\_touch} en el
firmware del nodo maestro.

% ---------------------------------------------------------------------------
%  5. RESULTADOS
% ---------------------------------------------------------------------------
\section{Resultados}
\label{sec:res}

\subsection{Configuración Global de Hardware}

La Tabla~\ref{tab:hw} resume la configuración de pines de todos los
nodos del sistema.

\begin{table}[H]
\centering
\caption{Configuración de hardware de los nodos}
\label{tab:hw}
\scriptsize
\begin{tabular}{p{1.2cm}p{1.4cm}p{1.4cm}p{1.6cm}}
\toprule
\textbf{Nodo} & \textbf{Interfaz} & \textbf{Pines} & \textbf{Periférico} \\
\midrule
Maestro & I2C & GPIO21/22 & DS3231 RTC \\
Maestro & SPI3 & 23/19/18/5 & MicroSD \\
Maestro & ADC1-CH6 & GPIO34 & Pulso (HR) \\
Nodo 2 & I2C & GPIO21/22 & MPU6050 \\
Nodo 3 & I2C & GPIO6/7 & QMI8658, CST816S \\
Nodo 3 & SPI2 & 11/10/9/8 & GC9A01 display \\
Nodo 3 & ADC & CH0 & Batería \\
\bottomrule
\end{tabular}
\end{table}

\subsection{Transmisión ESP-NOW y Tasa de Muestreo}

La Tabla~\ref{tab:perf} resume las tasas de muestreo y transmisión
configuradas en los nodos esclavos y la política de detección de
desconexión del maestro.

\begin{table}[H]
\centering
\caption{Rendimiento de comunicación del sistema}
\label{tab:perf}
\scriptsize
\begin{tabular}{p{1.9cm}p{1.3cm}p{1.5cm}p{1.1cm}}
\toprule
\textbf{Nodo} & \textbf{Muestreo} & \textbf{Envío ESP-NOW} & \textbf{Timeout} \\
\midrule
Nodo 2 (MPU6050) & 100 ms & 500 ms & 1000 ms \\
Nodo 3 (S3 Touch) & 50 ms & 50 ms & 1000 ms \\
Maestro (log SD) & -- & -- & 60 s/bloque \\
\bottomrule
\end{tabular}
\end{table}

El maestro detecta desconexión de cualquier nodo si no se recibe un
paquete válido en más de 1000~ms, activando las banderas
\texttt{mpu\_desconectado} y \texttt{touch\_desconectado} que se
reflejan en el log CSV.

\subsection{Gestión de Energía}

La Tabla~\ref{tab:power} muestra el impacto de las optimizaciones
energéticas implementadas.

\begin{table}[H]
\centering
\caption{Configuración de gestión de energía}
\label{tab:power}
\scriptsize
\begin{tabular}{p{3.0cm}p{3.7cm}}
\toprule
\textbf{Técnica} & \textbf{Configuración / Ahorro} \\
\midrule
DFS (maestro) & 240→80→10 MHz según carga \\
WiFi Power Save & WIFI\_PS\_MIN\_MODEM ($\approx$20 mA ahorro) \\
Light Sleep & Fallback si no soportado \\
TX power ESP-NOW & 10 dBm (40 unidades) \\
Pantalla (nodo 3) & Apagado auto tras 10 s sin táctil \\
\bottomrule
\end{tabular}
\end{table}

\subsection{Formato de Registro CSV}

Cada fila del archivo \texttt{datos.csv} en la tarjeta SD contiene un
registro temporal completo con granularidad de 1~segundo:

\begin{lstlisting}[style=csrc, numbers=none,
  caption={Cabecera del archivo datos.csv en la tarjeta SD},
  label=lst:csv]
timestamp,pulso_raw,
roll_mpu,pitch_mpu,ax_mpu,ay_mpu,az_mpu,
gx_mpu,gy_mpu,gz_mpu,
roll_tch,pitch_tch,...,
estado_mpu,estado_touch
\end{lstlisting}

Los campos \texttt{estado\_mpu} y \texttt{estado\_touch} indican si los
nodos esclavos estaban conectados en el instante del registro. El buffer
en RAM de 60 filas se vacía al disco cada 60~s, reduciendo el desgaste
de la tarjeta SD y el overhead de escritura en comparación con un volcado
fila a fila.

\subsection{Herramienta WHO\_AM\_I: Resultado Experimental}

La herramienta de diagnóstico (\texttt{herramienta\_who\_i\_am\_id\_ interno\_arduino\_ide.ino}) arrojó el siguiente resultado al ejecutarse
sobre el hardware del Nodo~2:

\begin{lstlisting}[style=csrc, numbers=none,
  caption={Salida del Monitor Serie de la herramienta diagnóstica},
  label=lst:out_who]
Leyendo WHO_AM_I...
WHO_AM_I: 0x72
\end{lstlisting}

La Tabla~\ref{tab:who} clasifica los posibles resultados.

\begin{table}[H]
\centering
\caption{Diagnóstico por valor del registro WHO\_AM\_I (0x75)}
\label{tab:who}
\scriptsize
\begin{tabular}{p{0.9cm}p{0.9cm}p{4.0cm}}
\toprule
\textbf{Valor} & \textbf{Estado} & \textbf{Acción recomendada} \\
\midrule
0x72 & \textcolor{green!60!black}{\bfseries OK} & Hardware del proyecto; driver acepta este valor \\
0x68 & \textcolor{green!60!black}{\bfseries OK} & MPU6050 original; sin cambios en driver \\
0x98 & \textcolor{orange}{\bfseries ALERTA} & Revisar máscara del driver \\
--- & \textcolor{red}{\bfseries ERROR} & Sin respuesta; verificar cableado \\
\bottomrule
\end{tabular}
\end{table}

% ---------------------------------------------------------------------------
%  6. DISCUSIÓN
% ---------------------------------------------------------------------------
\section{Discusión}
\label{sec:disc}

\subsection{Decisiones de Diseño Arquitectónico}

La elección de una topología estrella con ESP-NOW frente a alternativas
como BLE mesh o Zigbee se justifica por tres razones: (a)~los módulos
ESP32 ya incluyen el transceptor WiFi necesario sin costo adicional;
(b)~ESP-NOW opera en la misma banda de 2,4~GHz sin requerir un gateway
de red; y (c)~la implementación en ESP-IDF es directa y bien documentada
con ejemplos de código de referencia.

La separación de responsabilidades entre el Nodo~2 (movimiento de
tronco, MPU6050) y el Nodo~3 (interfaz táctil, QMI8658) permite una
modularidad que facilita el reemplazo o actualización de cualquier nodo
sin afectar los demás. La compatibilidad de la estructura
\texttt{datos\_imu\_t} entre nodos —garantizada por
\texttt{\_Static\_assert}— es un ejemplo de diseño defensivo que
previene errores silenciosos de compatibilidad de datos.

\subsection{Limitaciones del Sistema}

\begin{enumerate}
  \item \textbf{Sensor de pulso}: La lectura analógica por ADC es
        susceptible a ruido eléctrico y movimiento (artefactos de
        movimiento). Para uso clínico se recomendaría un sensor
        dedicado con interfaz I2C (p.ej. MAX30102).
  \item \textbf{Sincronización temporal}: La sincronización del RTC
        requiere intervención manual por consola. Un módulo NTP sobre
        WiFi proporcionaría sincronización automática.
  \item \textbf{Clasificación}: El sistema actual solo captura y
        registra datos; no implementa clasificación automática de
        AVD en tiempo real.
  \item \textbf{Alcance ESP-NOW}: El alcance típico en interiores es
        de 50--100~m en línea de visión directa, lo que puede
        ser insuficiente en entornos hospitalarios amplios.
\end{enumerate}

\subsection{Reproducibilidad y Código Abierto}

Todo el código fuente, incluyendo los drivers, las configuraciones de
CMake, los Doxyfiles de documentación y los binarios precompilados en
formato ELF y BIN, se encuentra disponible en el repositorio
GitHub~\cite{repo_github}. Cada carpeta de proyecto contiene su propio
\texttt{CMakeLists.txt} compatible con ESP-IDF~v6.0 y su
\texttt{devcontainer.json} para desarrollo en contenedores Docker,
facilitando la reproducción del entorno de compilación en cualquier
sistema operativo.

\subsection{Comparativa de Sensores IMU}

La Tabla~\ref{tab:imu_comp} compara las dos IMUs utilizadas en el sistema.

\begin{table}[H]
\centering
\caption{Comparativa entre MPU6050 y QMI8658}
\label{tab:imu_comp}
\scriptsize
\begin{tabular}{p{2.2cm}p{2.1cm}p{2.1cm}}
\toprule
\textbf{Parámetro} & \textbf{MPU6050} & \textbf{QMI8658} \\
\midrule
Ejes & 6-DOF & 6-DOF \\
Interfaz & I2C/SPI & I2C/SPI \\
Aceler. rango & $\pm$2--16~g & $\pm$2--8~g \\
Giro rango & $\pm$250--2000~°/s & $\pm$16--512~°/s \\
ADC resolución & 16 bits & 16 bits \\
Dirección I2C & 0x68/0x69 & 0x6A/0x6B \\
Nodo & Esclavo 2 & Esclavo 3 \\
\bottomrule
\end{tabular}
\end{table}

% ---------------------------------------------------------------------------
%  7. CONCLUSIONES Y TRABAJO FUTURO
% ---------------------------------------------------------------------------
\section{Conclusiones}
\label{sec:conc}

\begin{enumerate}
  \item Se diseñó e implementó un sistema wearable distribuido de tres
        nodos ESP32 para monitoreo de AVD, con comunicación inalámbrica
        por ESP-NOW, adquisición de datos IMU, frecuencia cardíaca, RTC
        y registro en tarjeta SD.

  \item La arquitectura maestro-esclavo con ESP-NOW proporciona
        latencias del orden de los milisegundos con bajo overhead de
        firmware, y la topología estrella simplifica la sincronización
        y el registro centralizado.

  \item El driver MPU6050 en C para ESP-IDF~v5+, con soporte de
        calibración y variantes de silicio (\texttt{WHO\_AM\_I} =
        \texttt{0x72}), constituye un componente reutilizable y
        documentado para proyectos de instrumentación inercial.

  \item El Nodo~3 demuestra la viabilidad de integrar IMU, pantalla
        circular táctil y comunicación inalámbrica en un único módulo
        ESP32-S3 de factor de forma compacto, con filtro complementario
        para estimación de ángulos de alta calidad.

  \item El código fuente completo, esquemas de conexión, binarios
        compilados y herramientas de diagnóstico están publicados en
        el repositorio GitHub~\cite{repo_github}, promoviendo la
        reproducibilidad y la transferencia tecnológica.
\end{enumerate}

\subsection{Trabajo Futuro}

\begin{itemize}
  \item Integrar algoritmos de clasificación de AVD (SVM, LSTM) en
        el nodo maestro aprovechando los 520~KB de SRAM del ESP32.
  \item Añadir transmisión a nube (MQTT/HTTPS) para monitoreo remoto
        en tiempo real.
  \item Implementar sincronización de tiempo NTP con fallback al RTC
        DS3231 para timestamping de alta precisión.
  \item Explorar reducción del consumo con ESP32-C6 y protocolo
        IEEE 802.15.4 (Thread/Zigbee) para redes de sensores corporales
        de mayor densidad.
\end{itemize}

% ---------------------------------------------------------------------------
%  REFERENCIAS
% ---------------------------------------------------------------------------
\begin{thebibliography}{9}

\bibitem{chen2012}
L. Chen, J. Hoey, C. D. Nugent, D. J. Cook y Z. Yu,
``Sensor-based activity recognition,'' \textit{IEEE Trans. Syst.,
Man, Cybern. Part C}, vol.~42, no.~6, pp.~790--808, Nov.~2012,
doi: 10.1109/TSMCC.2012.2198883.

\bibitem{anguita2012}
D. Anguita, A. Ghio, L. Oneto, X. Parra y J. L. Reyes-Ortiz,
``Human activity recognition on smartphones using a multiclass
hardware-friendly support vector machine,'' en \textit{Proc. IWAAL},
Vitoria-Gasteiz, 2012, pp.~216--223.

\bibitem{espressif2023}
Espressif Systems, \textit{ESP32 Technical Reference Manual}, v5.1,
Shanghai, 2023. [En línea]. Disp.:
\url{https://www.espressif.com/sites/default/files/documentation/esp32_technical_reference_manual_en.pdf}

\bibitem{espnow_docs}
Espressif Systems, \textit{ESP-NOW User Guide}, v2.0,
ESP-IDF Programming Guide, 2024. [En línea]. Disp.:
\url{https://docs.espressif.com/projects/esp-idf/en/latest/esp32/api-reference/network/esp_now.html}

\bibitem{invensense2013}
InvenSense Inc., \textit{MPU-6000 and MPU-6050 Product Specification},
Rev.~3.4, doc. PS-MPU-6000A-00, Sunnyvale, CA, 2013. [En línea]. Disp.:
\url{https://invensense.tdk.com/wp-content/uploads/2015/02/MPU-6000-Datasheet1.pdf}

\bibitem{nxp_i2c}
NXP Semiconductors, \textit{I2C-bus specification and user manual},
Rev.~7.0, UM10204, 2021. [En línea]. Disp.:
\url{https://www.nxp.com/docs/en/user-guide/UM10204.pdf}

\bibitem{kolban2017}
N. Kolban, \textit{Kolban's Book on ESP32}, Leanpub, 2017. [En línea].
Disp.: \url{https://leanpub.com/kolban-ESP32}

\bibitem{repo_github}
C. D. Cuellar Antury, ``Sistemas Embebidos — Actividades de la Vida
Diaria,'' Repositorio GitHub, Univ. del Cauca, 2024. [En línea]. Disp.:
\url{https://github.com/Cdcuellar/Sistemas-embebidos---Actividades-de-la-vida-diaria}

\end{thebibliography}

\end{document}
